\section{Nietzsche on the Will}

Nietzsche, who early on regarded Schopenhauer as one of his educators, came to reject his doctrine. The first step was to reject Schopenhauer's implicit supernaturalism, so for Nietzsche, the ``Will'' is part of the natural world and therefore has to be knowable, \emph{viz.}, a will is known by its purpose and Schopenhauer's Will is purposeless and therefore unintelligible. As a psychologist, Nietzsche looked within and recognized the Will as the ``Will to Power'', as the propensity to growth, creativity, expansion, strength, and so on as described in his analysis.

Since the Will to Power is the ultimate reality, it cannot be dependent on anything else. In particular, it cannot be dependent on Reason; quite to the contrary, Reason is a manifestation of the Will to Power and is of value only insofar as it enhances the Will. This means --- and I am speaking of the intellectual realm --- that there is no ultimate criterion of Truth to which the Will to Power must submit. Instead, there are competing perspectives, each vying for influence but none of them ultimately convincing.

There is a further step. The Will to Power is originally a positive concept, that is, it is descriptive of the actual state of affairs in the world. However, it becomes a normative, or moral, concept: what is stronger, that is, what has more power, becomes the Good. Nietzsche explains in \emph{The Antichrist}:

\begin{quotex}
What is good? --- All that heightens the feeling of power, the will to power, power itself in man. What is happiness? --- The feeling that power increases.
\end{quotex}


Note, however, that for Nietzsche the opposite of the Good is not Evil, but rather what is ``bad'', that is, the opposite of the ``feeling of power''. In particular, we can identify the Good with what is noble or strong or higher, generally what is more difficult and rare. The ``bad'', then, is whatever glorifies weakness.

In the battle of ideas, although there is neither ultimate truth nor a God to guarantee it, we can still measure strength in the force of an idea. We reject ideas that rely on blind belief, ignorance, or stupidity. Instead, we may choose to embrace ideas that are rational, factual, intricate, and uplifting. That is why Nietzsche seldom argues as a philosopher might, but rather relies on oracular pronouncements and insightful observations on individual and society.

Psychologically, we may question the reliance on the mere ``feeling'' of power. It is not uncommon to observe ``Nietzscheans'' puff themselves up to achieve that feeling, yet without being particularly powerful themselves. It seems there should be an objective standard to measure power, but then the Will to Power cannot be dependent on such a standard (or else the standard itself would be the ultimate reality.) There is also the problem that oftentimes Nietzscheans argue as if they were discerning ultimate truth and opposing ``Evil'' rather than the bad.

Metaphysically, we have to reject Nietzsche's naturalism and see the Will to Power as arising from something transcendent.


\flrightit{Posted on 2009-08-11 by Cologero}